\section{Обсуждение}

Наблюдаемая картина --- переоптимизация суррогатной награды~\cite{gao2022} с
конкретным, измеренным механизмом, а не общий «RL испортил модель».

\paragraph{Что произошло.} Политика не улучшила ранжирование, а нашла
область, где reward model систематически завышает оценку: короткие треки,
которые легче прослушать целиком, и небольшой набор айтемов, которые
ансамбль оценивает высоко при любом контексте. Обе стратегии увеличивают
$r_{\text{listen}}$, не требуя попадания в реальные цели пользователя.
Метрика NDCG на полном каталоге это показывает, а награда --- нет.

\paragraph{Почему защита не сработала.} Три причины, каждая проверяема по
артефактам:
\begin{enumerate}[leftmargin=1.2em,itemsep=1pt,topsep=2pt]
  \item ансамбль обучен на перекрывающихся подвыборках (80\% пользователей),
        поэтому разброс оценок мал и пессимистичный штраф
        $\bar r - 0.5\sigma_r$ смещает награду слабо;
  \item $\beta_{\text{KL}}=0.05$ к 64-му шагу держал KL на уровне 0.1 ---
        якорь разрешает существенный сдвиг, и этого хватает, чтобы им
        воспользоваться;
  \item награда однокомпонентна: оптимизируется только голова досмотра,
        поэтому падение головы «лайк» и коллапс разнообразия ничем не
        штрафуются.
\end{enumerate}

\paragraph{Почему reward hacking возможен здесь.} Голова досмотра различает
реальный исход пользователя от случайного трека (AUC 0.751), но на
правдоподобных кандидатах сигнал слаб (AUC «лайк» 0.564 на hard negatives,
Spearman досмотра 0.073). Политика оптимизирует градиент от неотличимых
кандидатов --- именно там оценка модели максимально уверена и максимально
неверна. Короткая длина трека и популярность --- свойства, которые модель
использует как прокси вместо пользовательского предпочтения.

\paragraph{Метрика разнообразия.} Gini по показам в топ-10 на 4\,600
пользователях вырождается в 1.0 у любой модели: почти каждый показанный
айтем встречается один раз на 630 тыс.\ каталога. Поэтому коллапс
измеряется межпользовательским Jaccard (0.030 $\to$ 0.736) и долей повторов
истории. Coverage@100 меняется слабо (3.4\% $\to$ 0.28\%), хотя выдача стала
почти идентичной: покрытие каталога --- плохой детектор коллапса, потому
что одинаковые списки у 4\,600 пользователей всё равно дают тысячи
уникальных айтемов.

\paragraph{Что это означает.} Рост награды и падение релевантности в этом
эксперименте не противоречат друг другу: они измеряют одну и ту же ошибку
с разных сторон. Награда --- внутренний сигнал модели, NDCG --- внешний.
Расхождение знаков голов reward model (досмотр растёт, лайк падает) делает
механизм явным: политика оптимизирует одну голову за счёт другой.
Долгосрочные цели (H5) в разомкнутом симуляторе не наблюдаемы по
построению и относятся к отдельной постановке с обратной связью.

\paragraph{Следующие шаги.} Не «ещё больше шагов RL», а изменение постановки:
многокомпонентная награда с явным штрафом за короткие треки и за совпадение
выдачи между пользователями, якорь $\beta_{\text{KL}} \ge 0.2$, обучение RM
на данных, порождённых политикой, отбор чекпоинта по NDCG на валидации, а
не по награде. Программа проверки --- отдельный прогон стадии
\code{ablate}.
